Esta pesquisa teve como objetivo investigar se diferentes representações computacionais (linguagens textuais, fluxogramas, \textit{blueprints} e plataformas \textit{Low-Code} e \textit{No-Code}) compartilham uma estrutura lógica comum capaz de servir de base para a tradução entre elas. Por meio de revisão bibliográfica e análise conceitual, esse objetivo foi alcançado: o trabalho identificou e caracterizou diferentes formas de representação e demonstrou, no plano teórico, a relação lógica que as une.

Verificou-se que, apesar de suas diferenças visuais e sintáticas, essas representações descrevem os mesmos elementos lógicos (decisões, condições, operações e sequências de execução), o que confirma a hipótese de que a lógica computacional pode ser compreendida de forma independente da forma que a expressa.

A principal contribuição do trabalho é o modelo conceitual proposto, no qual a lógica computacional atua como estrutura intermediária entre representações: uma representação de origem é interpretada, sua lógica é organizada nessa estrutura comum e, a partir dela, uma nova representação é gerada. O modelo reúne, sob uma mesma base teórica, representações que costumam ser tratadas como ambientes isolados.

Esse resultado tem implicações relevantes. Ao mostrar que código e representações visuais são formas distintas de uma mesma lógica, e não realidades incompatíveis, a pesquisa oferece uma base para repensar a interoperabilidade entre ambientes de desenvolvimento, favorecer a comunicação entre perfis técnicos distintos e aproximar os modelos textuais e visuais de construção de software.

Como continuidade natural, a validação prática do modelo e o desenvolvimento de mecanismos automáticos de conversão permitiriam avaliar empiricamente a abordagem e estendê-la a cenários mais complexos. Ainda assim, ao demonstrar a consistência teórica da separação entre lógica e representação, esta pesquisa oferece uma base sólida para futuras investigações sobre tradução e interoperabilidade entre representações computacionais.